%-----------------------------------------------------------------------
\section{Conclusion}
\label{sec:conclusion}
%-----------------------------------------------------------------------

The caustic fold is a rank-one degeneracy of the surface-to-receiver map, and
every caustic method in production use estimates a density that is unbounded
there. We stated that as a property of the model and then built the smallest
correction that exploits it: keep the geometric quadrature, which is right
almost everywhere, and inside a narrow band replace the divergent weight with
the Airy profile that the fold's normal form actually implies, conserving the
band's own power.

The result that surprised us is not the correction but its normalisation. The
Airy intensity is not integrable, so no constant normalises it
(Lemma~\ref{lem:notl1}), and the band's power is nevertheless finite
(Corollary~\ref{cor:integrable}). Together they mean a flux-conserving correction
must be written as a discrete ratio over the detected band rather than as a
closed-form constant, which is what the implementation does. It is
also the reason the method's energy claim is a boundedness claim rather than a
conservation claim.

On the comparative side, the position is narrower than a speed-up claim. Against
an independent brute-force trace we lead at coarse scales ($0.917$ against
$0.853$ SPPM, $0.828$ Cycles), tie SPPM at the finest scale measured
($0.721$ against $0.722$), and lose to plain geometric Monte Carlo below a
budget of about \num{262000} splats. What we do claim is that above roughly
\num{1000000} splats our single pass is $1.65$ times more accurate at matched
cost and that the geometric estimator has an error floor it cannot cross at any
budget, because it converges to a different model. That is a claim about
reachability, not about speed.

Two of our own earlier results are withdrawn here, both killed by estimators
rather than by physics: a dispersion model that was fitted rather than measured
and carried twice the true dispersion (Section~\ref{sec:verify-index}), and a
correlation score that was an artefact of a sign convention and a tile-naming
inversion (Section~\ref{sec:verify-metric}). Both are reported in the form that
makes them reproducible, because we think the failure mode is more likely to be
repeated than the method.

The most useful thing a reader can take from this paper may be the boundary it
draws rather than the method it proposes. Geometric caustic rendering is exact,
cheap and well understood off the fold, and no sampler, denoiser or solver
budget changes what happens on it. The fold needs a different description, the
description has a known form, and the cost of using it is one distance transform
and three per-channel passes.

\section{Reproduction}
\label{sec:repro}

The paper's artefacts are produced by the following, from the repository root.
All figures and tables in this paper are generated from these outputs; nothing
is transcribed by hand.

\begin{verbatim}
# the kinematic derivation's numerical constants
python tools/airy_facts.py                 # Airy peak, integral divergence
python tools/verify_refs.py                # citation verification (Crossref)

# wave-physics verification
python tools/verify_wave_correction.py     # profile vs scipy Airy
python tools/verify_ray_map.py             # the independent brute-force trace
python tools/fold_diag.py                  # fold locus, flux, boundary ablations
python tools/parity_cpp_vs_py.py           # C++ vs Python reference

# the comparison
python tools/spmm_vs_groundtruth.py        # three-way ranking, one axis
python tools/metric_parity_check.py        # both estimator conventions
python tools/accuracy_vs_ground_truth.py   # the headline accuracy
python tools/energy_audit.py               # Fresnel throughput
python tools/efficiency_study.py           # RMSE vs splat budget
python tools/ensemble_study.py             # 8-phase ensemble

# cost
build/bench.exe 512 1                      # bake timing, one thread
build/bench.exe 512 8                      # bake timing, eight threads

# the figures and clips
python tools/make_paper_figures.py         # paper figures
python tools/ripple_gifs.py                # the three ripple GIFs

# this paper
python tools/build_paper_pdf.py            # LaTeX -> PDF
\end{verbatim}

The three ripple animations referenced in the supplementary material are
produced by \verb|tools/ripple_gifs.py|, which renders the same animating
caustic through our estimator, Cycles and SPPM with one camera. The camera is
identical across the three to a projected-point disagreement of \num{0.0000}
(\verb|tools/cam_projection_check.py|), and the band uses the physical fringe
width rather than a visibility-exaggerated one, so the animations show the
measured dispersion rather than decorative fringes. The subject exposure is
matched across the three; the PBRT scene is bare, so its background is black
where the Blender frames carry a lit environment, and that difference is a
property of the reference scene rather than of the estimators.
